\section{Análise comparativa dos tipos societários}
\label{sec:comparacao}
A comparação considera composição, responsabilidade, abertura, manutenção e possibilidade de investimento. As expressões ``baixo'', ``moderado'' e ``elevado'' indicam complexidade relativa; não representam orçamento ou cotação de taxas (Lei nº 6.404/1976; Código Civil, 2002; LC nº 123/2006).

{\small
\begin{longtable}{@{}L{1.4cm} L{3.2cm} L{4.1cm} L{6.1cm}@{}}
\caption{Composição e responsabilidade das alternativas.}\label{tab:tipos}\\
\toprule
\cabfundo \cab{Forma} & \cab{Composição} & \cab{Responsabilidade} & \cab{Características e limites} \\
\midrule
\endfirsthead
\toprule
\cabfundo \cab{Forma} & \cab{Composição} & \cab{Responsabilidade} & \cab{Características e limites} \\
\midrule
\endhead
\bottomrule
\endfoot
MEI & Um titular; não admite sócios. & Não oferece a separação patrimonial típica da limitada; o titular responde pessoalmente pelas obrigações. & Regra geral: receita anual de até R\$~81 mil, proporcional no ano de abertura; até um empregado; apenas ocupações permitidas. A atividade efetiva deve constar da lista vigente. Não comporta os quatro fundadores. \\
\midrule
EI & Um empresário pessoa natural; não admite sócios. & Responsabilidade pessoal, sem limitação ao capital declarado, observadas as proteções legais de bens. & Não tem teto de receita inerente à forma jurídica; limites podem decorrer do porte ou regime fiscal. Estrutura simples, mas exposição patrimonial incompatível com a proposta conjunta. \\
\midrule
LTDA & Dois ou mais sócios nesta comparação; a lei também permite constituição unipessoal. & Restrita ao valor das quotas, com responsabilidade solidária dos sócios pela integralização do capital. & Contrato social, divisão de quotas e administração definida. Não há capital mínimo geral. Serviços não podem ser contribuição ao capital. Adequada a quatro fundadores. \\
\midrule
SLU & Um único sócio, pessoa natural ou jurídica. & Mesma regra de responsabilidade da sociedade limitada. & Sem capital mínimo geral. Pode admitir novos sócios mediante alteração contratual, permanecendo LTDA. Útil para fundador único, não para a composição proposta. \\
\midrule
S.A. & Em regra, dois ou mais subscritores; há exceções legais. & Limitada ao preço de emissão das ações subscritas ou adquiridas. & Capital dividido em ações; estatuto, órgãos e formalidades próprios. Companhia aberta depende de registro na CVM. Mais onerosa e complexa para a fase inicial. \\
\end{longtable}}

\clearpage
\subsection{Vantagens, limitações e adequação ao projeto}
\begin{description}[style=nextline,leftmargin=0cm]
\item[MEI.] A formalização simplificada e o recolhimento pelo DAS reduzem a complexidade para atividades permitidas. A vedação a sócios, o limite de receita e a necessidade de ocupação autorizada inviabilizam o modelo com quatro fundadores. Não se pressupõe que desenvolvimento de software ou intermediação digital sejam permitidos ao MEI.
\item[EI.] Facilita a organização de uma atividade individual, mas não admite participação dos demais integrantes e expõe o titular aos riscos de contratos, repasses e obrigações da operação.
\item[LTDA.] Permite distribuir quotas, definir administração e receber novos sócios, com governança proporcional à fase do negócio. Exige contabilidade, cumprimento contratual e controles financeiros. A limitação de responsabilidade não é proteção absoluta.
\item[SLU.] Concilia gestão individual e regime de responsabilidade da limitada. Escolhê-la apenas em nome de um integrante excluiria os demais da posição formal de sócio; por isso não é indicada sob a premissa de quatro fundadores.
\item[S.A.] Favorece estruturas de capital em ações e rodadas mais complexas. Em contrapartida, suas exigências de governança e manutenção são desproporcionais ao estágio acadêmico inicial. Uma S.A. não pode optar pelo Simples Nacional (LC nº 123/2006).
\end{description}

\subsection{Comparação operacional}
{\footnotesize
\renewcommand{\arraystretch}{1.05}
\begin{longtable}{@{}L{1.5cm} L{4.3cm} L{4.5cm} L{4.5cm}@{}}
\caption{Abertura, manutenção e investimento: avaliação qualitativa.}\\
\toprule
\cabfundo \cab{Forma} & \cab{Abertura e manutenção} & \cab{Entrada de investidores} & \cab{Conclusão para o PetCuida}\\
\midrule
\endfirsthead
\toprule
\cabfundo \cab{Forma} & \cab{Abertura e manutenção} & \cab{Entrada de investidores} & \cab{Conclusão para o PetCuida}\\
\midrule
\endhead
\bottomrule
\endfoot
MEI & Baixa complexidade, dentro das regras próprias. & Não admite sócios nem venda de quotas. & Inadequado à equipe e à estrutura projetada.\\
EI & Menor complexidade societária; obrigações fiscais conforme atividade. & Sem participação societária de terceiros nessa forma. & Inadequado à pluralidade de fundadores e ao risco.\\
LTDA & Moderada; contrato, registro e contabilidade. & Aumento de capital ou cessão de quotas; avaliar também instrumentos previstos em lei. & Melhor equilíbrio entre equipe, risco e expansão.\\
SLU & Similar à LTDA, com decisões do sócio único. & Admissão de sócios exige alteração contratual. & Alternativa somente se houver um fundador efetivo.\\
S.A. & Maior complexidade; possíveis simplificações legais não eliminam sua governança. & Emissão ou transferência de ações, conforme as regras aplicáveis. & Alternativa futura, não necessária agora.\\
\end{longtable}}